\section{A.1--A.2: stable parts and the integral Lefschetz problem}
\label{sec:A-foundations}

\subsection{Conventions and the witnesses}
We use the setting of Part A: $k$ is perfect of characteristic $p>0$,
$W=W(k)$, $X\subset\mathbf P^r_k$ is smooth, projective and connected of
dimension $n\geq2$, and $i:Y\hookrightarrow X$ is a smooth divisor in
$|\OO_X(d)|$, with $d\geq1$.  We put
\[
 C^j(Z)=H^j_{\cris}(Z/W),\qquad T^j(Z)=C^j(Z)_{\tors},\qquad
 E=\coker(C^{n-1}(X)\xrightarrow{i^*}C^{n-1}(Y)).
\]
All crystalline pullbacks are $W$-linear and commute with the
$\sigma$-semilinear crystalline Frobenius, where $\sigma$ is Witt
Frobenius.  On a finite-dimensional $k$-space with $\sigma$-semilinear
endomorphism $F$, the subscript $\ss$ means $\bigcap_{a\geq0}F^aM$.
It means the bijective part of $F$, not semisimplicity as a representation.

We have read the two supplied transcription witnesses for cote 4 and the
independent proof of 8.5.7.  The scans themselves were not supplied.
In particular, both witnesses retain a \emph{canceled} section 8.2 on
scans 3--4, including (8.2.1).  Thus the prompt's parenthesis that there is
``no section 8.2'' needs this qualification.  We make no independent
palaeographic assertion about the cancellation.

\begin{lemma}[Semilinear Fitting decomposition]\label{lem:A-fitting}
\status{Proved.} The functor $M\mapsto M_{\ss}$ on finite-dimensional
$k$-spaces with $\sigma$-semilinear endomorphism is exact.  There is a
functorial decomposition into a bijective summand and a nilpotent summand.
The analogous stable-part functor is exact on finite-length $W$-modules.
For finite $W$-modules we write $M^{\rm u}$ for its $p$-adic limit; this
functor is exact, commutes with reduction modulo $p$, and selects the
slope-zero part after inverting $p$, together with any bijective torsion.
\end{lemma}
\begin{proof}
Since $\sigma$ is an automorphism, images and kernels of iterates are
linear submodules.  For a finite-length module both sequences stabilize.
For large $a$, the map $F$ is bijective on $\im F^a$, nilpotent on
$\ker F^a$, and their intersection is zero.  Their lengths sum to the
length of the original module, proving the decomposition.  A commuting
map preserves both summands.  In a short exact sequence its restriction
to bijective summands is surjective: apply a sufficiently large common
iterate of $F$ to a lift of a vector in the bijective quotient, and use
bijectivity there.  This also proves exactness in the middle.

Apply the finite-length construction to $M/p^aM$ and take the inverse
limit.  The transition maps on bijective summands are surjective by the
exactness just proved.  Completeness gives the asserted summand of $M$.
To check exactness for a submodule $N\subset M$, use the induced
filtration $N\cap p^aM$.  The Artin--Rees lemma makes this filtration
cofinal with the $p$-adic filtration of $N$; inverse limits of the
resulting exact finite-length sequences are exact by the surjectivity
of transition maps.  Reduction modulo $p$ follows either from the
decomposition or from the sequence for multiplication by $p$.
On the complement, some power of $F$ sends $M$ into $pM$, so all its
rational slopes are positive.  On the bijective summand all slopes are
zero.  The finite-length assertion and this limit formulation agree
with \citet[Lemma 4.8 and Remark 4.10]{Ogus1975}.
\end{proof}

\begin{lemma}[Recheck of R1, with no lower degree bound]\label{lem:A-R1}
\status{Proved.} Frobenius acts nilpotently on
$H^j(X,\OO_X(-d))$ for $j<n$.  Consequently
\[
 H^j(X,\OO_X)_{\ss}\longrightarrow H^j(Y,\OO_Y)_{\ss}
\]
is an isomorphism for $j<n-1$ and is injective for $j=n-1$.
\end{lemma}
\begin{proof}
The ideal $I=\OO_X(-d)$ is preserved by taking $p$th powers, and its
$a$th Frobenius iterate factors as the $\sigma^a$-semilinear map
$I\to I^{p^a}=\OO_X(-p^ad)$ followed by the inclusion into $I$.
For $j<n$, Serre duality identifies the dual of
$H^j(X,\OO_X(-p^ad))$ with
$H^{n-j}(X,\omega_X(p^ad))$, which vanishes for large $a$ by Serre
vanishing.  Thus an iterate is zero on $H^j(X,I)$.  Apply
Lemma~\ref{lem:A-fitting} to the cohomology sequence of
$0\to I\to\OO_X\to i_*\OO_Y\to0$.
This confirms the precise factorization used in the supplied independent
proof; it requires no Kodaira vanishing and no assumption that the full
coherent restriction is injective.
\end{proof}

\begin{lemma}[Recheck of R2]\label{lem:A-R2}
\status{Proved.} For every smooth proper $k$-scheme $Z$, the edge map
\[
 \alpha:H^j_{\dR}(Z/k)\longrightarrow H^j(Z,\OO_Z)
\]
induces an isomorphism on bijective Frobenius parts.
\end{lemma}
\begin{proof}
Let $\alpha:\Omega^\bullet_{Z/k}\to\OO_Z[0]$ be augmentation and
let $\beta:\OO_Z[0]\to\Omega^\bullet_{Z/k}$ be the
$\sigma$-semilinear chain map $f\mapsto f^p$ in degree zero.
The latter is a chain map because $d(f^p)=0$.  As maps of complexes of
abelian sheaves, with the indicated scalar twists, absolute Frobenius
on the de Rham complex is $\beta\alpha$, and Frobenius on $\OO_Z$ is
$\alpha\beta$.  On the bijective parts the inverse of $\alpha$ is
$\beta F^{-1}$, interpreted with the inverse scalar twist.  This proves
the claim without degeneration of the Hodge-to-de Rham spectral
sequence.  It is the actual argument needed from R2.
\end{proof}

\begin{lemma}[Integral Nakayama form]\label{lem:A-derived-Nakayama}
\status{Proved.} Let $K$ be a bounded complex with finite cohomology
over a complete discrete valuation ring $R$ of uniformizer $p$.
If $K\otimes_R^{\mathbf L}R/p\in D^{\geq s}$, then
$K\in D^{\geq s}$ and $H^s(K)$ is $p$-torsion-free.
\end{lemma}
\begin{proof}
The universal coefficient sequence is
\[
0\to H^j(K)/p\to H^j(K\otimes_R^{\mathbf L}R/p)
  \to H^{j+1}(K)[p]\to0.
\]
For $j<s$, Nakayama gives $H^j(K)=0$.  Taking $j=s-1$ then gives
$H^s(K)[p]=0$.  Over a discrete valuation ring this is equivalent to
torsion-freeness for a finite module.
\end{proof}

\begin{theorem}[An explicit sufficient integral Lefschetz bound]\label{thm:A-high-degree}
\status{Proved under stated hypotheses.} Write $\operatorname{reg}$
for Castelnuovo--Mumford regularity with respect to the specified
embedding, and set
\[
 d_{\rm int}(X)=\max\bigl(1,\operatorname{reg}(\Omega_X^j):0\leq j\leq n\bigr).
\]
For every $d\geq d_{\rm int}(X)$ and every smooth
$Y\in|\OO_X(d)|$, crystalline restriction is an isomorphism in degrees
$j<n-1$ and is injective with torsion-free cokernel in degree $n-1$.
In particular, A.1(a)--(c) hold in this range, and restriction induces
$T^{n-1}(X)\simeq T^{n-1}(Y)$.
The full coherent restriction in degree $n-1$ is injective whenever
$d\geq\max(1,\operatorname{reg}\omega_X)$.
\end{theorem}
\begin{proof}
For $t\geq d_{\rm int}(X)$, regularity gives
$H^b(X,\Omega_X^a(t))=0$ for $b>0$ and all $a$.
By the wedge pairing and Serre duality,
\[
 R\Gamma(X,\Omega_X^a(-t))\in D^{\geq n}.
\]
This holds in particular for $t=md$ and every $m\geq1$.
For $0\leq a\leq n-1$, resolve the exterior powers of the conormal
sequence
\[
0\to\OO_Y(-d)\to\Omega_X^1|_Y\to\Omega_Y^1\to0.
\]
The terms preceding $\Omega_X^a|_Y$ are
$\Omega_X^{a-b}|_Y(-bd)$, $1\leq b\leq a$.
Their cohomology lies in $D^{\geq n-1}$ by the restriction sequence and
the preceding vanishing.  The cone of
$R\Gamma(\Omega_X^a)\to R\Gamma(\Omega_X^a|_Y)$ lies in
$D^{\geq n-1}$.  Resolving the additional kernel with the indicated
$a$ terms proves
\[
\operatorname{Cone}\bigl(R\Gamma(\Omega_X^a)
 \to R\Gamma(\Omega_Y^a)\bigr)\in D^{\geq n-a-1}.
\]
For $a=n$, the same bound is immediate from $\Omega_Y^n=0$.
After the de Rham shift $[-a]$, all these cones lie in
$D^{\geq n-1}$.  The finite filtration by form degree therefore gives
the same bound for the cone of the de Rham restriction.

Crystalline cohomology is a perfect $W$-complex here, and its derived
reduction modulo $p$ is the de Rham complex.  Apply
Lemma~\ref{lem:A-derived-Nakayama} to its restriction cone.  The long
cohomology sequence now proves the stated isomorphisms and the
injectivity; the degree-$(n-1)$ cokernel is a submodule of the
torsion-free group $H^{n-1}$ of the cone.  If $A\hookrightarrow B$ is
an injection of finite $W$-modules with torsion-free quotient, then its
torsion subgroups agree: every torsion element of $B$ maps to zero in
the quotient.  This proves the torsion assertion.
This supplies an explicit regularity bound in the argument of
\citet[Lemma 2.12]{BMS2016}, which we have checked in the supplied
source.  Finally, Serre duality identifies
$H^{n-1}(X,\OO_X(-d))^\vee$ with $H^1(X,\omega_X(d))$; regularity
and the divisor exact sequence prove the last assertion.
\end{proof}

\begin{theorem}[All-degree integral Lefschetz in two further ranges]
\label{thm:A-all-degree-ranges}
\status{Proved under stated hypotheses.}
All the conclusions of Theorem~\ref{thm:A-high-degree}, including
full coherent injectivity, hold for every $d\geq1$ in either of the
following two cases:
\begin{enumerate}[label=(\roman*)]
\item $X\subset\mathbf P^r_k$ is a smooth global complete intersection;
\item $n\leq p$ and $X$ admits a smooth lift over $W_2(k)$.
\end{enumerate}
More generally, in case (ii) with no dimension restriction, crystalline
restriction is an isomorphism below degree $\min(n,p)-1$, and is
injective with torsion-free cokernel in that degree.  Its torsion
subgroups therefore agree in that degree as well.
\end{theorem}
\begin{proof}
We first give the exact input from \citet[Definition 2.2.1 and
Theorem 2.2.3]{CotnerZavyalov2024}.  For a smooth ample divisor
$D\subset Z$, where $Z$ lifts to $W_2(k)$, their theorem puts the
cone on degree-$a$ Hodge cohomology in
$D^{\geq\min(\dim Z,p)-1-a}$.  For a smooth complete intersection
$Z\subset\mathbf P^r$, the corresponding cone from projective
space lies in $D^{\geq\dim Z-a}$.  We read the proof, including the
conormal-exterior-power induction and its precise characteristic
bound.  This does not require a lift of the particular divisor.

After shifting the Hodge terms by $[-a]$, the finite de Rham
filtration gives the same bound without $-a$ for the de Rham cone.
Derived crystalline base change and
Lemma~\ref{lem:A-derived-Nakayama} give the integral assertion in
case (ii).  Setting $a=0$ gives coherent injectivity in the required
degree when $n\leq p$.

In case (i), both $X$ and $Y$ are smooth global complete
intersections in the same projective space: adjoining the defining
equation of the section to the regular sequence for $X$ gives the
regular sequence for $Y$.  Apply the projective-space assertion
separately to $X$ and $Y$.  The cone of
$R\Gamma_{\dR}(\mathbf P^r/k)\to R\Gamma_{\dR}(X/k)$ lies in
$D^{\geq n}$, while the analogous cone for $Y$ lies in
$D^{\geq n-1}$.  The triangle of these two cones identifies the
cone of $X\to Y$ as an extension of the second cone by the first
cone shifted by $[1]$.  It therefore lies in $D^{\geq n-1}$.
The same integral Nakayama argument proves crystalline injectivity,
saturation, and the lower-degree isomorphisms.  The identical
cone argument for $a=0$ proves coherent injectivity.  Finally,
an injective map with torsion-free cokernel identifies the torsion
subgroups, as proved in Theorem~\ref{thm:A-high-degree}.
\end{proof}

\begin{theorem}[The stable exact sequence of 8.5 in every degree]\label{thm:A-exact85}
\status{Proved.} No additional lower bound on $d$ is necessary for
the exact sequence
\begin{equation}\label{eq:A-exact85}
0\longrightarrow T^n(X)[p]_{\ss}
 \longrightarrow (E/pE)_{\ss}
 \xrightarrow{\varphi} E(Y,\OO_Y)_{\ss}
 \longrightarrow0.
\end{equation}
In particular, $\ker v=0$, and (8.5.5)--(8.5.6) become
\begin{align}
\dim_k(E/pE)_{\ss}
 &=\dim_k E(Y,\OO_Y)_{\ss}+\dim_k T^n(X)[p]_{\ss},\label{eq:A-dimension}\\
\dim_k(E/pE)_{\ss}&\geq\dim_k E(Y,\OO_Y)_{\ss}.
\end{align}
Both directions of 8.5.10 hold: $\varphi$ is an isomorphism if and
only if $T^n(X)[p]_{\ss}=0$.
\end{theorem}
\begin{proof}
Derived crystalline base change and Lemma~\ref{lem:A-R2} turn the
universal coefficient sequence into a natural short exact sequence
\begin{equation}\label{eq:A-stable-UCT}
0\to(C^j(Z)/p)_{\ss}\to H^j(Z,\OO_Z)_{\ss}
 \to C^{j+1}(Z)[p]_{\ss}\to0.
\end{equation}
If $j>\dim Z$, its middle term vanishes, so
$(C^j(Z)/p)_{\ss}=0$.  This also implies $C^j(Z)[p]_{\ss}=0$.
Indeed some iterate of $F$ maps $C^j(Z)$ into $pC^j(Z)$; repeated
iteration annihilates its torsion, as is seen after any $W$-module
splitting into free and torsion modules.  Thus $C^n(Y)[p]_{\ss}=0$.

Use \eqref{eq:A-stable-UCT} in degree $j=n-1$ for $X$ and $Y$ and
the vertical restriction maps.  The middle vertical arrow is injective
by Lemma~\ref{lem:A-R1}.  Hence the left vertical arrow is injective
as well, and the lower row has zero right-hand term.  The snake lemma
gives
\[
0\to C^n(X)[p]_{\ss}\to
 \coker\bigl((C^{n-1}(X)/p)_{\ss}\to(C^{n-1}(Y)/p)_{\ss}\bigr)
 \to\coker v\to0.
\]
Right exactness of reduction modulo $p$ and exactness of stable
extraction identify the middle term with $(E/pE)_{\ss}$ and the last
term with $E(Y,\OO_Y)_{\ss}$.  The map is the edge-map-induced
$\varphi$ of the manuscript.  Taking dimensions proves the formulas
and exactness proves the equivalence.
\end{proof}

\begin{theorem}[Unit-root count without full integral saturation]\label{thm:A-unit-count}
\status{Proved.} Suppose $k=\mathbf F_q$.  For every $d\geq1$ as above,
the group $E^{\rm u}$ is torsion-free and
$\dim_k(E/pE)_{\ss}$ equals the number, with multiplicity, of
$p$-adic unit reciprocal roots of
\[
 P_{\rm van}(T)=
 \frac{\det(1-T\operatorname{Frob}_q^{\rm geom}\mid
 H^{n-1}_{\et}(Y_{\bar k},\mathbf Q_\ell))}
 {\det(1-T\operatorname{Frob}_q^{\rm geom}\mid
 H^{n-1}_{\et}(X_{\bar k},\mathbf Q_\ell))},\qquad\ell\ne p.
\]
The assertion is independent of $\ell$ and of an embedding of its
algebraic reciprocal roots into $\overline{\mathbf Q}_p$.
\end{theorem}
\begin{proof}
We first work over $\bar k$.  The Artin--Schreier sequence and
semilinear Lang's theorem identify
$H^j_{\et}(Z,\mathbf F_p)\otimes\bar k$ with
$H^j(Z,\OO_Z)_{\ss}$; the map $F-1$ is surjective on every finite
dimensional Frobenius module over $\bar k$, including its nilpotent
part.  Lemma~\ref{lem:A-R1} thus gives mod-$p$ weak Lefschetz.
Apply Lemma~\ref{lem:A-derived-Nakayama} over $\mathbf Z_p$ to
\[
K=\operatorname{Cone}\bigl(R\Gamma_{\et}(X_{\bar k},\Zp)
 \longrightarrow R\Gamma_{\et}(Y_{\bar k},\Zp)\bigr).
\]
Here continuous cohomology is derived $p$-complete; finite mod-$p$
cohomology, its uniform vanishing outside a bounded range, and the
coefficient sequences give the required finite cohomology.  Alternatively
this finiteness follows at once from the integral comparison below.
The cone modulo $p$ lies in $D^{\geq n-1}$.  Consequently the integral
map in degree $n-1$ is injective and its cokernel is torsion-free.

\citet[Theorem 4.9, Remark 4.10 and Corollary 4.11]{Ogus1975}
identify, functorially and \emph{including torsion},
\[
 C^j(Z)^{\rm u}\simeq
 H^j_{\et}(Z_{\bar k},\Zp)\otimes_{\Zp}W(\bar k).
\]
The embedding hypothesis in the cited theorem holds: a projective
$k$-scheme embeds into projective space over every $W_a(k)$ by
composition with its special-fibre inclusion.  Stable extraction is
exact by Lemma~\ref{lem:A-fitting}; thus $E^{\rm u}$ is the scalar
extension of the preceding torsion-free cokernel.  Faithfully flat
descent proves its torsion-freeness over $W(k)$.  Its rank is the
difference of the slope-zero dimensions for $Y$ and $X$.

By \citet[Theorem 1, p.~74]{KatzMessing1974}, the characteristic
polynomial in each degree is the same in rational crystalline and
$\ell$-adic cohomology.  The ordinary $\ell$-adic weak Lefschetz theorem
makes the displayed quotient the polynomial on the vanishing quotient.
Therefore its number of unit reciprocal roots is precisely the
difference just obtained.  Since $(E/pE)_{\ss}=E^{\rm u}/pE^{\rm u}$,
this is the required dimension.  The equality is an equality of
polynomials in $\mathbf Z[T]$, with algebraic roots; any chosen
$p$-adic realization gives the crystalline slope multiset and the
same count.  No algebraicity conjecture for a Lefschetz inverse is used.

Our Frobenius convention is worth spelling out.  Crystalline
$F^{[k:\Fp]}$ has the same characteristic polynomial as geometric
Frobenius on $\ell$-adic cohomology.  On the $F=1$ local system the
arithmetic generator acts by its inverse.  This agrees with the
insertion ``l'inverse du'' in the witness to 4:7.
\end{proof}

\begin{proposition}[The transpose in 8.2.1]\label{prop:A-transpose-u2}
\status{Proved under the stated identification.} The transpose of
$i^*:T^n(X)[p]\to T^n(Y)[p]$ under crystalline torsion duality is
\[
i_*:T^{n-1}(Y)/p\longrightarrow T^{n+1}(X)/p.
\]
If $i^*:T^{n-1}(X)\to T^{n-1}(Y)$ is an isomorphism, as it is in
Theorem~\ref{thm:A-high-degree}, this becomes
\[
 u_2'':T^{n-1}(X)/p\longrightarrow T^{n+1}(X)/p,
 \qquad x\longmapsto d\,c_1^{\cris}(\OO_X(1))\smile x.
\]
For a fixed degree its kernel is independent of the particular smooth
section.  Across different degrees it need not be unchanged: the map
is zero when $p\mid d$ and has the same kernel as cup product with
$c_1^{\cris}(\OO_X(1))$ when $p\nmid d$.
\end{proposition}
\begin{proof}
Derived Poincare duality gives perfect linking pairings
\[
 T^j(Z)\times T^{2\dim Z+1-j}(Z)\longrightarrow W[1/p]/W.
\]
Pairing the $p$-torsion of the first factor with reduction modulo $p$
of the second yields a perfect $k$-pairing.  Pullback and Gysin are
adjoint, by the projection formula for the trace.  Setting $j=n$ for
$X$ and $Y$ gives exactly the displayed degrees, without shifting
either torsion index.  Under the indicated identification, the
self-intersection formula gives $i_*i^*(x)=c_1(\OO_X(Y))\smile x$,
and $\OO_X(Y)=\OO_X(d)$.  Reduction modulo $p$ proves the final
statements.  This interpretation does not itself prove the torsion
identification at degrees below $d_{\rm int}(X)$.
\end{proof}

\begin{proposition}[The two sequences on the sheet 8.4.4]\label{prop:A-star-dual}
\status{Proved with explicit operator convention.} Define $V$ to be
the $\sigma^{-1}$-semilinear Poincare adjoint of $F$ on the
complementary cohomological degree.  On rational cohomology of $X$,
this is Frobenius pushforward and satisfies $VF=FV=p^n$.
Superscripts ${\ss}$ below refer to this $V$, not to the usual
degree-independent Verschiebung of a Dieudonne module.  Then the
sequences transcribed as $(*)$ and $(**)$ are dual:
\begin{align*}
0&\to(T^n(X)/p)_{\ss}\to H^n(X,\OO_X)_{\ss}
 \to(C^n(X)_{\rm free}/p)_{\ss}\to0,\\
0&\to(C^n(X)_{\rm free}/p)^{\ss}\to
 H^0(X,\omega_X)^{\ss}\to T^{n+1}(X)[p]^{\ss}\to0.
\end{align*}
\end{proposition}
\begin{proof}
The term $C^{n+1}(X)[p]_{\ss}$ vanishes by
\eqref{eq:A-stable-UCT} and the argument following it.  Hence
$H^n(X,\OO_X)_{\ss}\simeq(C^n(X)/p)_{\ss}$, and the first sequence
is the stable part of the torsion--free quotient sequence.
Dualize it by $\Hom_k(-,k)$.  Serre duality identifies the dual of
$H^n(\OO_X)$ with $H^0(\omega_X)$; free crystalline Poincare duality
identifies the dual of the free middle-degree module with itself;
the linking pairing identifies the dual of $T^n/p$ with
$T^{n+1}[p]$.  These identifications send the transpose of $F$ to
the specified $V$, and preserve the bijective parts.  They give the
second sequence with precisely its reversed arrows.  In particular,
one must not transpose an $F$-stable torsion group and leave the
operator labelled $F$: the dual assertion uses $V$.
\end{proof}

\begin{question}[The remaining small-degree problem in A.1]\label{open:A-small-degree}
\status{Open in this working dossier.} Outside the two all-degree
ranges of Theorem~\ref{thm:A-all-degree-ranges}, for
$1\leq d<d_{\rm int}(X)$,
we have not determined in complete generality whether $E$ is
torsion-free, whether $i^*$ induces the asserted torsion isomorphism,
or whether it is injective on full crystalline cohomology.  We have
also not decided full coherent injectivity in this range.  A failure
of $H^{n-1}(X,\OO_X(-d))$ to vanish does not by itself refute
injectivity: its map to $H^{n-1}(X,\OO_X)$ must be nonzero.  In
dimension two, if $Y$ is geometrically connected, that map is
injective, so a smooth very ample Kodaira counterexample would
suffice.  The published examples inspected so far have not yet been
verified to satisfy all these extra conditions.
\end{question}
